\section{Conclusion}

This project explored reinforcement learning for continuous control and 
sim-to-real transfer across two tasks. In Part 1, the Actor-Critic algorithm 
equipped with Generalized Advantage Estimation, advantage normalization, 
and learning-rate decay proved to be the strongest performer on the Hopper 
environment, achieving average rewards around 2000 and single episodes 
exceeding 3000, while vanilla REINFORCE and its constant-baseline variant 
served as important stepping stones toward understanding the role of 
variance reduction in policy gradient methods.

In Part 2, the focus shifted to sim-to-real transfer in the PandaPush environment. The experimental results demonstrated
a clear advantage of Soft Actor-Critic (SAC) over Proximal Policy Optimization (PPO), which failed to solve the task within
the available training budget. Among the transfer strategies deployed to fill the reality gap,
Uniform Domain Randomization (UDR) achieved the best overall performance, outperforming both vanilla SAC and
unexpectedly Automatic Domain Randomization (ADR) too.

The final selected approach is therefore SAC combined with UDR. By training on a broad distribution of cube masses,
UDR produced a substantially sturdier policy and reduced the sim-to-real gap by approximately 66\% relative to
the non-randomized baseline. These results confirm the effectiveness of domain randomization as a practical strategy
for improving transfer performance in robotic reinforcement-learning tasks.